\section{Introduction}
\label{sec:intro}

\subsection{Background and motivation}
Skin cancer is among the most common cancers worldwide, and early recognition of malignant lesions such as
melanoma and basal cell carcinoma strongly improves outcomes. Dermatoscopy (dermoscopy) is a non-invasive
imaging technique that magnifies sub-surface skin structures, but reading dermatoscopic images reliably takes
years of training. Deep convolutional neural networks (CNNs) have reached dermatologist-level accuracy on this
problem when trained on very large image collections~\cite{esteva2017}. Most clinical settings, however, have only
a few thousand labelled images, and these are heavily skewed towards benign nevi. In that low-data, imbalanced
regime, \emph{transfer learning} is the standard remedy: a network pretrained on a large generic dataset
(ImageNet) is reused, and only adapted to the target domain~\cite{yosinski2014,kornblith2019}.

This assignment follows \textbf{Task~2 (CNN / Transfer Learning / Fine-tuning)} of the CT-1 brief:
\emph{fine-tune a pretrained CNN (e.g.\ ResNet, EfficientNet) on a domain-specific dataset (e.g.\ medical
images)}, with exposure to \emph{transfer learning, freezing/unfreezing layers and domain adaptation}. We chose
medical images: the seven-class DermaMNIST benchmark~\cite{medmnistv2}, derived from the HAM10000 dermatoscopic
collection~\cite{ham10000}. We use the 224$\times$224 version released with MedMNIST+~\cite{medmnistplus}, so
ImageNet models see images at their native input resolution.

\subsection{Problem statement}
Given an RGB dermatoscopic image $x \in [0,255]^{224\times224\times3}$, predict one of seven diagnostic
categories $y \in \{\texttt{akiec},\texttt{bcc},\texttt{bkl},\texttt{df},\texttt{mel},\texttt{nv},\texttt{vasc}\}$.
A useful model must do well on the \emph{rare} classes, including melanoma (\texttt{mel}). A model that always
predicts the majority class \texttt{nv} already reaches \MajorityAcc\% test accuracy, so plain accuracy alone
can mislead. We therefore treat \textbf{macro-F1} as the primary metric and also report balanced accuracy,
macro one-vs-rest AUC and calibration.

\subsection{Objectives}
\begin{enumerate}
  \item Analyse the dataset: class balance, colour statistics and possible train/test leakage.
  \item Build a reproducible PyTorch pipeline: preprocessing, augmentation, imbalance-aware training,
        validation-based model selection and evaluation.
  \item Compare transfer-learning strategies on two pretrained backbones (ResNet-50~\cite{resnet} and
        EfficientNet-B0~\cite{efficientnet}) against a CNN trained from scratch:
        (i)~\emph{feature extraction} (frozen backbone), (ii)~\emph{full fine-tuning} with discriminative learning
        rates, and (iii)~\emph{gradual unfreezing}~\cite{ulmfit}.
  \item Measure the effect of class-weighted loss (ablation) and of test-time augmentation (TTA).
  \item Explain the predictions with Grad-CAM~\cite{gradcam} and t-SNE embeddings, and discuss limitations.
\end{enumerate}

\subsection{Contributions of this report}
\begin{itemize}
  \item A controlled comparison of \NumRuns{} training configurations. All of them use the same data pipeline,
        the same seed and the same validation-based checkpoint selection; the test set is used only once, for
        the final numbers.
  \item An honest evaluation: bootstrap 95\% confidence intervals, per-class metrics, calibration (ECE) and a
        near-duplicate audit of the official splits.
  \item A public, documented codebase that rebuilds every number and figure in this report with one command
        (Section~\ref{sec:codebase}).
\end{itemize}
